%%%%%%%%%%%%%%%%%%%%%%%%%%%%%%%%%%%%%%%%%%%%%%%%%%%%%%%%%%%%%%%%%%%%%%
%                                                                    %
% Mean Ces\`aro-type summability  of Fourier-Neumann series,         %
% by                                                                 %
% \'Oscar Ciaurri, Krzysztof Stempak and Juan L. Varona              %
%                                                                    %
% DATE: first version, 29-October-2002                               %
% DATE: last modification, 23-October-2003                           %
%                                                                    %
% Written with AMS-LaTeX 2.0                                         %
%                                                                    %
%%%%%%%%%%%%%%%%%%%%%%%%%%%%%%%%%%%%%%%%%%%%%%%%%%%%%%%%%%%%%%%%%%%%%%

\documentclass[a4paper]{amsart}
%\usepackage{amsmath,amsthm} %%% amsthm necessary to use ``proof''

% \usepackage[notcite,notref]{showkeys}
% \overfullrule=5pt

%\date{October 23, 2003}
% \newcommand{\datename}{\textit{Date}:} %% <- Original
\renewcommand{\datename}{THIS PAPER HAS BEEN PUBLISHED IN:}
\date{\textit{Studia Sci.\@ Math.\@ Hungar.\@} \textbf{42} (2005), 413--430.}

%%%%%%%% A LOT OF BIG SEVERAL-LINE EQUATIONS (aligns, etc.).
%%%%%%%% THUS, WE ALLOW TO BREAK THEM IN PAGES.
\allowdisplaybreaks

%%%% \span already exists in TeX
\newcommand{\spa}{\operatornamewithlimits{span}}
\newcommand{\supp}{\operatornamewithlimits{supp}}
\newcommand{\iden}{\operatorname{Id}}

\newcommand{\calM}{\mathcal{M}}
\newcommand{\calR}{\mathcal{R}}
\newcommand{\calS}{\mathcal{S}}
\newcommand{\calH}{\mathcal{H}}
\newcommand{\calJ}{\mathcal{J}}
\newcommand{\mua}{d\mu_\alpha}
\newcommand{\muag}{d\mu_{\alpha+\gamma}}

\newtheorem{theor}{Theorem}
\newtheorem{propo}{Proposition}
\newtheorem{corol}[theor]{Corollary}

%%%%%%%%%%%%%%%%%%%%%%%%%%%%%%%%%%%%%%%%%%%%
\begin{document}
%%%%%%%%%%%%%%%%%%%%%%%%%%%%%%%%%%%%%%%%%%%%

%%%%%%%%%%%%%%%%%%%%%%%%%%%%%%%%%%%%%%%%%%%%%%
\title[Mean Ces\`aro-type summability of Fourier-Neumann series]{Mean
Ces\`aro-type summability of Fourier-Neumann series}

%\author{\'Oscar Ciaurri, Krzysztof Stempak and Juan L. Varona}

\author[ \'O. Ciaurri]{\'Oscar Ciaurri}
\address{Departamento de Matem\'aticas y Computaci\'on,
Universidad de La Rioja,
Edificio J.~L.~Vives,
Calle Luis de Ulloa s/n,
26004 Logro\~no, Spain}
\email{oscar.ciaurri@dmc.unirioja.es}

\thanks{Research of the first and third authors supported by grant
BFM2003-06335-C03-03 of the DGI} %% BFM2003-06335-C03-03

\author[K. Stempak]{Krzysztof Stempak}
\address{Instytut Matematyki,
Politechnika Wroc\l{a}wska,
Wyb.~Wyspia\'nskiego 27,
50-370 Wro\-c\l{a}w,
Poland}
\email{stempak@im.pwr.wroc.pl}

\thanks{Research of the second author supported in part by
KBN grant\#~2 P03A 028 25}

\author[J. L. Varona]{Juan L. Varona}

\address{Departamento de Matem\'aticas y Computaci\'on,
Universidad de La Rioja, Edificio J.~L.~Vives,
Calle Luis de Ulloa s/n, 26004 Logro\~no, Spain}
\email{jvarona@dmc.unirioja.es}

% \subjclass{Primary 42C10; Secondary 44A20}
\subjclass[2000]{Primary 42C10; Secondary 44A20}

\keywords{Riesz summation process, Ces\`aro means, summability of
Fourier-Neumann series, Hankel transform, Bessel functions}

%%%%%%%%%%%%%%%%%%%%%%%%%%%%%%%%%%%%%%%%%%%%%%
\begin{abstract}
Let $J_\nu$ be the Bessel function of order $\nu$. For $\alpha>-1$, the
functions $x^{-\alpha-1}J_{\alpha+2n+1}(x)$, $n=0,1,2\dots$, form an
orthogonal system in $L^2(x^{2\alpha+1}\,dx)$, but the span of such functions
is not dense in this space. For a function $f$, let $S_k^{\alpha}f$ denote
the $k$th partial sum of the Fourier-Neumann series of $f$. In this paper we
provide the minimal conditions on a real $\gamma$
and $1<p<\infty$, for which the means
$R_n^{\alpha}f = \frac{\lambda_0 S_0^{\alpha}f + \cdots
+ \lambda_n S_n^{\alpha}f}{\lambda_0 + \cdots + \lambda_n}$,
$\lambda_k = 2(\alpha+2k+2$), are uniformly bounded in the spaces
$L^p(x^{2(\alpha+\gamma)+1}\,dx)$. Clearly, the convergence
$R_n^{\alpha}f\to f$ holds only for functions from the closure
of the linear span of the orthogonal system in these spaces.
As a byproduct of the main result, we obtain a characterization of the closure of
the span in terms of functions whose modified Hankel transforms of order $\alpha$
are supported on the interval $[0,1]$.
\end{abstract}
%%%%%%%%%%%%%%%%%%%%%%%%%%%%%%%%%%%%%%%%%%%%%%

%%%%%%%%%%%%%%%%%%%%%%%%%%%%%%%%%%%%%%%%%%%%%%
\maketitle
%%%%%%%%%%%%%%%%%%%%%%%%%%%%%%%%%%%%%%%%%%%%%%

%%%%%%%%%%%%%%%%%%%%%%%%%%%%%%%%%%%%%%%%%%%%%%
\section{Introduction and statement of results}
%\setcounter{equation}{0}
\label{sec:intro}
%%%%%%%%%%%%%%%%%%%%%%%%%%%%%%%%%%%%%%%%%%%%%%

Let $J_\nu$ stand for the Bessel function of the first kind of order~$\nu$.
For $\alpha>-1$, the formula
\begin{equation*}
  \int_0^\infty J_{\alpha+2n+1}(x)J_{\alpha+2m+1}(x)\frac{dx}{x}
  = \frac{\delta_{nm}}{2(2n+\alpha+1)},
  \quad n,m = 0,1,2,\dots,
\end{equation*}
provides an orthonormal system $\{j_n^\alpha\}_{n=0}^\infty$ in
$L^2((0,\infty),x^{2\alpha+1}\,dx)$ ($L^2(\mua)$ and, more generally, $L^p(\mua)$
from now on) given by
\begin{equation*}
  j_n^\alpha(x) = \sqrt{2(2n+\alpha+1)}\,J_{\alpha+2n+1}(x) x^{-\alpha-1},
  \qquad n = 0,1,2,\dots\,.
\end{equation*}
For a function $f$, provided that the coefficients
\begin{equation*}
  c_k^\alpha(f) = \int_0^\infty f(y)j_k^\alpha(y) y^{2\alpha+1}\,dy,
  \qquad k = 0,1,2,\dots,
\end{equation*}
exist, consider its formal expansion
$f\sim \sum_{k=0}^\infty c_k^\alpha(f) j_k^\alpha(x)$ and the partial sum
operators
\begin{equation*}
  S_n^\alpha(f,x) = \sum_{k=0}^n c_k^\alpha(f) j_k^\alpha(x),
  \qquad n = 0,1,2,\dots\,.
\end{equation*}
Series of the form $\sum_{n\ge0} a_n J_{\alpha+n}$ are usually
called the Neumann series, hence we refer to $\sum_{k=0}^\infty c_k^\alpha(f)
j_k^\alpha(x)$ as to the Fourier-Neumann series.

For $\alpha$ and $\gamma$, $\alpha>-1$, $\gamma>-1-\alpha$, let
\begin{equation*}
  p_0(\alpha,\gamma) = \max\left\{1,\frac{4(\alpha+\gamma+1)}{2\alpha+3}\right\},
  \qquad p_1(\alpha,\gamma) =
  \begin{cases}
    \frac{4(\alpha+\gamma+1)}{2\alpha+1},&\text{ $\alpha>-1/2$}, \\
    \infty, &\text{$-1<\alpha\le-1/2$}.
  \end{cases}
\end{equation*}
In the case $\gamma=0$ we simply write $p_i(\alpha)$ in place of $p_i(\alpha,0)$, $i=1,2$.

The essential aim of this paper is to study the convergence $R_n^\alpha f\to f$
in the $L^p(\muag)$ spaces (note that it is equivalent to the
study of the convergence in the weighted $L^p(x^{2\gamma}\mua)$ spaces),
$1<p<\infty$, where $R_n^\alpha$ are means of the partial sum operators
$\{S_k^\alpha\}_{k=0}^n$, cf.~\eqref{eq:Rna}.
To pose the problem correctly, we must impose some natural conditions on the
parameters involved. First, we have the requirement $j_n^\alpha\in L^p(\muag)$.
Second, to ensure the existence of the coefficients $c_n^\alpha(f)$ for every
$f \in L^p({\muag})$, we must require
$j_n^\alpha \in L^{p'}(x^{-2\gamma p'}\,\muag)$.
However, assuming $1<p<\infty$, we have
$j_n^\alpha\in L^p(\muag)\cap L^{p'}(x^{-2\gamma p'}\,\muag)$
for every $n=0,1,2,\dots$ if and only if
$0 < \alpha+\gamma+1 < (\alpha+1)p$ and
$p_0(\alpha,\gamma) < p < p_1(\alpha,\gamma)$,
and these are precisely the minimal assumptions we impose.
The last equivalence easily follows by using the following
well-known asymptotics for the Bessel functions:
\begin{equation}
\label{eq:zero}
  J_\nu(x) = \frac{x^\nu}{2^\nu \Gamma(\nu+1)} + O(x^{\nu+2}),
    \qquad x \to 0^+;
\end{equation}
\begin{equation}
\label{eq:infty}
  J_\nu(x)=\left(\frac{2}{\pi x}\right)^{1/2}\left[ \cos\left(x-\frac{\nu\pi}2
    - \frac\pi4 \right) + O(x^{-1}) \right], \qquad x \to \infty.
\end{equation}

Most orthogonal systems considered in the literature, for instance, the
trigonometric, Jacobi, Laguerre and Hermite systems as well as the ones
related to Freud weights, Bessel and Dini, are complete. Moreover, the
$\spa$ of each of these systems is dense in the corresponding $L^p$ spaces,
$1\le p<\infty$.
Thus, the $L^p$-convergence of the associated Fourier series is equivalent
to the uniform boundedness of the partial sum operators.
However, this is not the present case since the span of $\{j^\alpha_n\}$ is not
dense in~$L^p(\mua)$.
Therefore, it is important to identify the closure in $L^p(\mua)$
of the subspace spanned by $j_n^\alpha$, $n = 0,1,2,\dots$,
or, more generally, the space
\begin{equation*}
  B_{p,\alpha,\gamma} = \overline{\spa \{j_n^\alpha : n = 0,1,2,\dots\}}
  \qquad (\text{closure in } L^p(\muag)),
\end{equation*}
because the uniform boundedness of the partial sum operators $S_n^\alpha$,
or the Ces\`aro-type means $R_n^\alpha$, would imply the $L^p$-convergence
only for functions in $B_{p,\alpha,\gamma}$. Note
that the spaces $B_{p,\alpha,\gamma}$ are well defined if
$j_n^\alpha\in L^p(\muag)$ and this holds if and only if
$\alpha+\gamma+1>0$ and $p_0(\alpha,\gamma)<p$ (again by using~\eqref{eq:zero}
and~\eqref{eq:infty}, or see~\cite{Ci}).

Identifying the spaces $B_{p,\alpha,\gamma}$ with those defined in
\eqref{eq:Epa}, we will use the modified Hankel transform $H_{\alpha}$ given by
\begin{equation}
\label{eq:htr}
  H_{\alpha}(f,x) = \int_0^\infty \frac{J_\alpha(xy)}{(xy)^\alpha} \, f(y)
    y^{2\alpha+1} \, dy,
  \qquad x>0.
\end{equation}
For $\alpha\ge-1/2$, \eqref{eq:zero}, \eqref{eq:infty}
and H\"older's inequality show that~\eqref{eq:htr} is well defined for every
$f\in L^p(\mua)$ with $1 \le p <p_0(\alpha)$.
Furthermore, it is easy to see that $H_{\alpha}$ is a bounded
operator from $L^1(\mua)$ into $L^\infty(\mua)$.
It is also well known that $H_{\alpha}$ extends to an isometric
isomorphism from $L^2(\mua)$ onto itself and $H_{\alpha} \circ H_{\alpha} =
\iden$ (for $-1<\alpha<-1/2$,
the isometric extension of $H_{\alpha}$ to $L^2(\mua)$ can
also be done; see~\cite{Dur,BeSt}).
As a consequence, by interpolation, for $\alpha\ge-1/2$ we consider $H_{\alpha}$ to
be a bounded operator from $L^p(\mua)$ into $L^{p'}(\mua)$, $1\le p\le2$.

The Hankel transform of the function $j_n^\alpha$ is
\begin{equation*}
   H_{\alpha}(j_n^\alpha,x) = \sqrt{2(\alpha+2n+1)} P_n^{(\alpha,0)}(1-2x^2)
   \chi_{[0,1]}(x),
\end{equation*}
where $P_n^{(\alpha,\beta)}(x)$ is the $n$-th Jacobi polynomial of order
$(\alpha,\beta)$, cf.~\cite{CiGPV}, a remark on p.~127.
The important fact here is that $H_{\alpha}(j_n^\alpha)$ is supported on
$[0,1]$; then, only functions $f$ having this property could be approximated
by the sequence $S_n^\alpha f$. We describe this property using the language
of the modified Hankel transform.

Thus, define the operator $M_{\alpha}$, associated with $H_{\alpha}$, as the multiplier
operator
\begin{equation*}
  M_{\alpha} f = H_{\alpha}(\chi_{[0,1]}H_{\alpha} f).
\end{equation*}
Applying Fubini's theorem, we have
$M_\alpha (f,x) = \int_0^\infty M_\alpha(x,y)f(y)\,d\mu_\alpha(y)$,
where
\begin{equation}
\label{eq:Makernel}
  M_\alpha(x,y) = (xy)^{-\alpha} \int_0^1 J_\alpha(xs)J_\alpha(ys) s\,ds.
\end{equation}
It is known that $M_\alpha$ is a bounded operator on $L^p(\mua)$
if and only if $p_0(\alpha)<p<p_1(\alpha)$, cf.~\cite{Her,V,CiGPV-DIE,CiV,BeSt}.

For $\alpha \ge -1/2$ and $p_0(\alpha) < p < p_1(\alpha)$ define the
following closed subspace of~$L^p(\mua)$:
\begin{equation}
\label{eq:Epa}
  E_{p,\alpha} = \{ f \in L^p(\mua): M_\alpha f = f\}
  = M_\alpha(L^p(\mua)).
\end{equation}
The spaces $E_{p,\alpha}$ enjoy some nice properties: for $1<p<q<\infty$,
$E_{p,\alpha} \subset E_{q,\alpha}$ and the inclusion is continuous and dense.
Besides, the dual space of $E_{p,\alpha}$ is naturally identified
with $E_{p',\alpha}$, cf.~\cite{V} for details.

Note that we assume $\alpha>-1$ in the definition of $B_{p,\alpha,\gamma}$.
 We require, however, $\alpha\ge-1/2$ for $E_{p,\alpha}$. Actually, the
definition of $E_{p,\alpha}$ can be extended to the whole range of $\alpha >
-1$. But in the case $\alpha < -1/2$, $H_{\alpha}$ is not a bounded operator
from $L^1$ into $L^\infty$ since $J_\alpha(t)t^{-\alpha}$ is not a bounded function.
As a consequence, the spaces
$E_{p,\alpha}$ do not behave for $\alpha<-1/2$ like those for $\alpha\ge-1/2$.
Thus, some of the results in this paper will be established for $\alpha>-1$,
but we will require $\alpha\ge-1/2$ when $E_{p,\alpha}$ appears.

It is clear from the very definition that for $f \in E_{p,\alpha}
\cap L^2(\mua)$ the Hankel
transform of $f$ is supported on $[0,1]$. When $p=2$, using the fact that
$H_{\alpha}$ is an isomorphism
and the completeness of the Jacobi system, we easily get that
$\{j_n^\alpha\}$ is, indeed, a basis for~$E_{2,\alpha}$.
But this cannot be proved in such a direct way for other values of~$p$.

In~\cite{V}, one of the authors
(actually, in that paper there is a change of notation that is, basically,
a change of variable $x \mapsto x^2$; see \cite[Remark~1]{CiV} for
appropriate comments on the change of notation) proved the following:
by using the basic case $p=2$, and the properties
of $E_{p,\alpha}$, it follows that $B_{p,\alpha,0} = E_{p,\alpha}$
for $2 \le p < p_1(\alpha)$.
On the other hand, it was also shown that, for $\alpha\ge-1/2$,
\begin{equation*}
  \|S_n^\alpha f\|_{L^p(\mua)} \le C \|f\|_{L^p(\mua)},\qquad n=0,1,2,\dots,
\end{equation*}
if and only if $\max\{p_0(\alpha),4/3\} < p < \min\{p_1(\alpha),4\}$
(for a more general result in a weighted setting, see \cite[Theorem~1]{CiGPV}).
Then, for this range of $p$'s, $\|S_n^\alpha f-f\|_{L^p(\mua)} \to 0$
for every $f$ from the space $B_{p,\alpha,0}$, already identified
as $E_{p,\alpha}$ for $2 \le p < p_1(\alpha)$.
By using duality and the uniform boundedness of the partial sum operators $S_n^\alpha$,
it  also follows that  $B_{p,\alpha,0} = E_{p,\alpha}$
when $\max\{p_0(\alpha),4/3\} < p < 2$.

If $-1/2 \le \alpha < 0$, we have $ p_0(\alpha)<4/3$. Thus, there is a gap
$p_0(\alpha) < p < 4/3$, for which  the previous argument
does not show if $B_{p,\alpha,0}$ and $E_{p,\alpha}$ are equal.
Filling this gap is another  aim of this paper (see
Theorem~\ref{th:Bp=Ep} with $\gamma=0$).

To overcome the difficulty, we will use means equivalent to Ces\`aro
$(C,1)$ means instead of partial sums. This will allow us to
eliminate the condition $4/3 < p < 4$ and,
again by using a duality argument, will enable us to
identify $B_{p,\alpha,0}$ with $E_{p,\alpha}$ in the whole range of
$p_0(\alpha) < p < p_1(\alpha)$.
Actually, instead of using Ces\`aro means, we will use the means
\begin{equation}
\label{eq:Rna}
  R_n^\alpha =
  \frac{\lambda_0 S_0^\alpha + \cdots + \lambda_n S_n^\alpha}
    {\lambda_0 + \cdots + \lambda_n},
  \qquad \lambda_k = 2(\alpha+2k+2).
\end{equation}
A reason for this is that for such means it is possible to write an
exact expression for their kernels.
This summation method is sometimes known as the Riesz method
$(R,\lambda_k)$; see, for instance, \cite[\S\,59, p.~470]{Kno}.
The convergence in this method is equivalent
to the convergence of Ces\`aro $(C,1)$ means
(see, for instance, \cite[\S\,3.8, Th.~14]{Har}).
Thus, we refer to this method as to a \textit{Ces\`aro-type} summation method.

Investigating the uniform boundedness and convergence of
$R_n^\alpha$ reveals that the kernels of these operators remind
those for the Hankel transform
partial sum operators, cf.~\cite{CiV}; boundedness properties of the
corresponding operators were studied there by using a
convolution structure for the Hankel transform. Unfortunately, in
the case of kernels considered here, an extra factor of size
$xy/n^2$ appears, and the method used in~\cite{CiV} didn't turn out to be
useful. On the other hand, however, the kernels have a similar
behaviour as the ones that appear in~\cite{CiStV}; boundedness
properties of the corresponding operators were established without
using properties of the aforementioned convolution. It occurs that
the proofs may be adapted to the present case and the extra factor
can be properly annihilated. This is particularly well visible in the proof
of Proposition~\ref{pro:BcAc}.

Finally, some comments seem to be relevant.
Although, in many aspects, the situation of the present setting is
similar to the one of the Laguerre expansions, for instance, the
condition $4/3<p<4$ disappears when we pass from the partial sum
operators to the $(C,1)$ means, cf.~\cite{Po} (see~\cite{St03} for a
motivation of such a comparison), there are important differences. First,
the $(C,1)$ (and also $(C,\delta)$) summability of the Laguerre series
can be deduced from an associated convolution structure,
see~\cite{St91}. Second, we always restrict the $p$'s interval to
$1<p<\infty$, not including $p=1$ or $p=\infty$. This is explained by the fact that
already in the unweigted case, $\gamma=0$, the system $j_n^\alpha$, $n=0,1,2,\dots$,
is not contained in $L^1(d\mu_\alpha)$.

The main result of the paper states the uniform boundedness of
$R_n^\alpha$ in $L^p(\muag)$.

%%%%%%%%%%%%%%%%%%%%%%%%%%%%%%%%%
\begin{theor}
\label{th:main}
Let $\alpha > -1$ and $\gamma\in\mathbb{R}$. Assume that
\begin{equation}
\label{eq:condition1}
  0 < \alpha+\gamma+1 < (\alpha+1)p
\end{equation}
and
\begin{equation}
\label{eq:condition2}
  p_0(\alpha,\gamma) < p < p_1(\alpha,\gamma).
\end{equation}
Then
\begin{equation*}
  \|R_n^\alpha f\|_{L^p(\muag)} \le C \|f\|_{L^p(\muag)}\,,
\end{equation*}
with $C$ independent of~$n=0,1,2,\dots$ and~$f$.
\end{theor}

%%%%%%%%%%%%%%%%%%%%%%%%%%%%%%%%%
Note, that the assumptions that appear in Theorem~\ref{th:main} are,
according to what was said earlier, ``natural'' (minimal).

The second main result of the paper concerns the spaces~$B_{p,\alpha,\gamma}$.

%%%%%%%%%%%%%%%%%%%%%%%%%%%%%%%%%
\begin{theor}
\label{th:Bp=Ep} Let $\alpha > -1$, $\alpha+\gamma\ge-1/2$,
$p_0(\alpha+\gamma)<p<p_1(\alpha+\gamma)$ and $\gamma<1/2$.
If $p<2$, assume further that
$p_0(\alpha,\gamma)<p$
and $\alpha+\gamma+1<\frac{p}4(2\alpha+4\gamma+3)$.
Then $B_{p,\alpha,\gamma} = E_{p,\alpha+\gamma}$.
\end{theor}
%%%%%%%%%%%%%%%%%%%%%%%%%%%%%%%%%

The proofs of Theorems~\ref{th:main} and~\ref{th:Bp=Ep} are postponed to
Sections~\ref{sec:proofmain} and~\ref{sec:Bp=Ep}. Theorem~\ref{th:Bp=Ep} imediately
shows that $B_{p,\alpha,0}=E_{p,\alpha}$ in the whole range of $p$'s in
which both spaces are well defined.

By using Theorems~\ref{th:main} and~\ref{th:Bp=Ep}, the proof of the
following result is standard.

%%%%%%%%%%%%%%%%%%%%%%%%%%%%%%%%%%%%%%
\begin{corol}
Let $\alpha > -1$, $\alpha+\gamma\ge-1/2$,
$p_0(\alpha+\gamma)<p<p_1(\alpha+\gamma)$, $\gamma<1/2$,
and~\eqref{eq:condition1} and
\eqref{eq:condition2} are satisfied. If
$p<2$, assume further $\alpha+\gamma+1<\frac{p}4(2\alpha+4\gamma+3)$.
Then, $R_n^\alpha f\to f$ in
$L^{p}(\muag)$ for every $f\in E_{p,\alpha+\gamma}$.
\end{corol}
%%%%%%%%%%%%%%%%%%%%%%%%%%%%%%%%%%%%%%%

Apart of the expansions based on the system $\{j_n^\alpha\}_{n=0}^\infty$, there
is another setting in which the Fourier-Neumann series may be  investigated.
The system of functions
\begin{equation*}
 \varphi_n^\alpha(x) = \sqrt{2(2n+\alpha+1)}\,J_{\alpha+2n+1}(x) x^{-1/2},
  \qquad n = 0,1,2,\dots\,
\end{equation*}
is orthonormal (but not complete) in
$L^2((0,\infty), dx)$; in what follows we simply write $L^2(dx)$ or, more generally,
$L^p(x^\gamma dx)$ for weighted $L^p$ spaces.
Hence, for any suitable function $f$ on $(0,\infty)$ consider its expansion
$f\sim\sum_{k=0}^\infty d_k^\alpha(f) \varphi_k^\alpha(x)$, $d_k^\alpha(f)=
\int_0^\infty f(x)\varphi_k^\alpha(x)\,dx$, and the partial sum operators
\begin{equation*}
  \calS_n^\alpha(f,x) = \sum_{k=0}^n d_k^\alpha(f) \varphi_k^\alpha(x),
  \qquad n = 0,1,2,\dots\,.
\end{equation*}
Analogously, define the means
\begin{equation*}
\label{eq:Rnacal}
  \calR_n^\alpha =
  \frac{\lambda_0 \calS_0^\alpha + \cdots + \lambda_n \calS_n^\alpha}
    {\lambda_0 + \cdots + \lambda_n},
  \qquad \lambda_k = 2(\alpha+2k+2).
\end{equation*}
Clearly enough, results concerning both types of expansions are linked together since
$\calS_n^\alpha(f,x) =x^{\alpha+1/2}S_n^\alpha((\cdot)^{-(\alpha+1/2)}f,x)$.

Therefore, Theorem~\ref{th:main} has its counterpart in the following.

%%%%%%%%%%%%%%%%%%%%%%%%%%%%%%%%%
\begin{theor}
\label{th:maincal}
Let $\alpha > -1$, $\gamma\in\mathbb{R}$ and $1<p<\infty$. Assume that
\begin{equation}
\label{eq:cal}
  \max\{-1,-1-p(\alpha+1/2)\} < \gamma < \min\{-1+p, -1+p(\alpha+3/2)\}.
\end{equation}
Then
\begin{equation*}
  \|\calR_n^\alpha f\|_{L^p(x^\gamma\,dx)} \le C \|f\|_{L^p(x^\gamma\,dx)}\,,
\end{equation*}
with $C$ independent of~$n=0,1,2,\dots$ and~$f$.
\end{theor}
%%%%%%%%%%%%%%%%%%%%%%%%%%%%%%%%%

As in the case of Theorem~\ref{th:main}, it can be easily seen by a direct computation
that the assumptions of Theorem~\ref{th:maincal} are minimal in the sense that
for all $n=0,1,\dots$, $\varphi_n^\alpha\in L^p(x^\gamma\,dx)$ and the coefficients
$ d_k^\alpha(f)$ exist for every $f\in L^p(x^\gamma\,dx)$ if and
only if~\eqref{eq:cal} is satisfied.
Clearly, Theorem~\ref{th:maincal} and Theorem~\ref{th:main} are equivalent.
Actually, the direct proof of Theorem~\ref{th:maincal} simplifies a little bit;
for instance, the formula corresponding
to that of Proposition~\ref{pro:3sum} becomes
$$
  \calR_n^\alpha(f,x) = \calM_{\alpha}^0(f,x)
  - \frac{d_n}{2} \Big(x\calM_{\alpha+1}^1(yf,x)-x
    \calM_{\alpha+2n+3}^1(yf,x)\Big).
$$

%%%%%%%%%%%%%%%%%%%%%%%%%%%%%%%%%%%%%%%%%%%%%%
\section{The kernel}
%\setcounter{equation}{0}
\label{sec:kernel}
%%%%%%%%%%%%%%%%%%%%%%%%%%%%%%%%%%%%%%%%%%%%%%

In this section we find exact expressions for the kernel of the Ces\`aro-type
means~$R^\alpha_n$.
The partial sum operators $S_n^\alpha $ can be written as
\begin{equation}
\label{eq:Sn}
  S^\alpha_n(f,x) = \int_0^\infty S_n^\alpha(x,y) f(y) y^{2\alpha+1}\,dy
\end{equation}
with the kernel
\begin{equation*}
  S_n^\alpha(x,y) = (xy)^{-(\alpha+1)}
    \sum_{k=0}^n 2(\alpha+2k+1) J_{\alpha+2k+1}(x)J_{\alpha+2k+1}(y).
\end{equation*}
Thus, by~\eqref{eq:Rna},
\begin{equation*}
  R^\alpha_n(f,x) = \int_0^\infty R_n^\alpha(x,y)f(y)y^{2\alpha+1}\,dy
\end{equation*}
with
\begin{equation*}
  R_n^\alpha(x,y)
  = d_n \left(\lambda_0 S_0^\alpha(x,y) + \cdots
  + \lambda_n S_n^\alpha(x,y)\right),
  \qquad \lambda_k = 2(\alpha+2k+2),
\end{equation*}
where $d_n^{-1} = \sum_{k=0}^n \lambda_k=4(n+1)(\alpha+n+2)$.

For $\nu>-1$, let
\begin{equation}
\label{eq:phi1}
  \calJ_\nu(x,y) = \int_0^1 J_\nu(xs) J_\nu(ys) s\,ds.
\end{equation}
Then
\begin{equation}
\label{eq:phi2}
  \calJ_\nu(x,y) = \frac{1}{x^2-y^2}\{ xJ_{\nu+1}(x)J_{\nu}(y)
    - yJ_{\nu}(x)J_{\nu+1}(y) \}
\end{equation}
and
\begin{equation}
\label{eq:phi3}
  \calJ_\nu(x,y) = \frac{1}{x^2-y^2}\{ yJ'_{\nu}(y)J_{\nu}(x)
  - xJ'_{\nu}(x)J_{\nu}(y) \}.
\end{equation}
We applied von Lommel's formula, cf.~\cite{CiGPV}, and, for the last equality,
the identity $zJ_{\nu+1}(z) = \nu J_{\nu}(z) - zJ'_{\nu}(z)$.

By using~\eqref{eq:phi2}, it can be proved (see~\cite{CiGPV}) that
\begin{equation}
\label{eq:kernelKn}
  S_n^\alpha(x,y) = (xy)^{-\alpha} \{\calJ_\alpha(x,y)
  - \calJ_{\alpha+2n+2}(x,y)\}.
\end{equation}
Then
\begin{multline*}
  R_n^\alpha(x,y) = \frac{1}{2}d_n (xy)^{-\alpha}
    \sum_{k=0}^n 2(\alpha+2k+2) \{\calJ_\alpha(x,y)
      - \calJ_{\alpha+2k+2}(x,y)\} \\
  \begin{aligned}
  &= (xy)^{-\alpha} \left[\calJ_\alpha(x,y)
    - d_n \int_0^1 s\bigg( \sum_{k=0}^n 2(\alpha+2k+2)
    J_{\alpha+2k+2}(xs) J_{\alpha+2k+2}(ys)\bigg) \,ds\right] \\
  &= (xy)^{-\alpha}\left[\calJ_\alpha(x,y)
    - d_n \int_0^1 s (xsys) (\calJ_{\alpha+1}(xs,ys)
    - \calJ_{\alpha+2n+3}(xs,ys)) \,ds\right].
  \end{aligned}
\end{multline*}
Therefore, denoting
\begin{equation*}
  \xi_\nu(x,y) = 2 \int_0^1 s^3 \calJ_\nu(xs,ys) \,ds,
\end{equation*}
we finally obtain
\begin{equation}
\label{eq:Jn3sum}
  R_n^\alpha(x,y) = (xy)^{-\alpha} \calJ_\alpha(x,y)
  - \frac{1}{2}d_n(xy)^{-\alpha+1}( \xi_{\alpha+1}(x,y)
  - \xi_{\alpha+2n+3}(x,y)).
\end{equation}
In this decomposition, the first summand is just the integral kernel of the
multiplier operator $M_\alpha$, as one can observe by comparing~\eqref{eq:phi1}
with~\eqref{eq:Makernel}.
Its boundedness properties has already been studied;
see for instance~\cite{CiGPV}
(note also that this summand appears in the decomposition
of the kernel $S_n^\alpha(x,y)$ used to study the uniform
boundedness of~$S_n^\alpha$).

We will find another expression for $\xi_\nu(x,y)$. By using~\eqref{eq:phi1},
the change of variable $t=w/s$ and Fubini's theorem easily produce
\begin{equation*}
  \xi_\nu(x,y)=\int_0^1 w(1-w^2) J_\nu(xw) J_\nu(yw) \,dw.
\end{equation*}
This expression for $\xi_\nu(x,y)$ will be useful for our purposes for the
following reasons: let $\calH_\nu$ denote the (non-modified) Hankel
transform of order~$\nu > -1$, i.e.,
\begin{equation*}
  \calH_\nu(f,x)
  = \int_0^\infty (xy)^{1/2}J_\nu(xy)f(y)\,dy, \qquad x>0.
\end{equation*}

The Bochner-Riesz multiplier of order $\delta\ge0$ for
$\calH_\nu$ is given by
\begin{equation}
\label{eq:B-Rd}
  \calM_{\nu}^\delta(f,x)
  = \calH_\nu(m^\delta \cdot \calH_\nu f,x)
  = \int_0^\infty \calM_{\nu}^\delta(x,y)f(y)\,dy,
\end{equation}
where $m^\delta(y) = (1-y^2)^\delta_+$ and
\begin{equation}
\label{eq:KB-Rd}
  \calM^{\delta}_{\nu}(x,y)
  = \int_0^1 (1-w^2)^\delta (xw)^{1/2}J_\nu(xw)(yw)^{1/2}J_\nu(yw)\, dw.
\end{equation}
Then, $\xi_\nu(x,y) = (xy)^{-1/2}\calM^{1}_{\nu}(x,y)$.
Pointwise estimates for $\calM^{1}_{\nu}(x,y)$,
as well as necessary and sufficient conditions for the uniform boundedness
\begin{equation*}
  \|\calM_{\nu}^1(f,x)x^a\|_{L^p(dx)}
  \le C\|f(x)x^a\|_{L^p(dx)},
  \qquad \nu\ge\alpha,
\end{equation*}
were found in~\cite{CiStV}. We will partially use those results
in Section~\ref{sec:proofmain} of this paper.

Actually, also the first summand in~\eqref{eq:Jn3sum} can be described
following the notation of~\eqref{eq:B-Rd}:
the multiplier $M_\alpha$ (for the modified Hankel transform $H_{\alpha}$)
is related to the  multiplier $\calM_{\alpha} = \calM_{\alpha}^0$
(for the Hankel transform $\calH_\alpha$) by mean of
$\calM_\alpha(f,x) = x^{\alpha+1/2}M_{\alpha}((\cdot)^{-(\alpha+1/2)}f,x)$,
with the analogous relation for the kernels
$\calM_\alpha(x,y)$ and $M_{\alpha}(x,y)$.

Summarizing, we found the following decomposition.

%%%%%%%%%%%%%%%%%%%%%%%%%%%%%%%%%
\begin{propo}
\label{pro:3sum}
Let $\alpha>-1$ and $n=0,1,2\dots\,$. Then
\begin{multline*}
  R_n^\alpha(f,x)
  = x^{-\alpha-1/2} \calM_{\alpha}^0(y^{\alpha+1/2}f,x) \\
  - \frac{d_n}{2} \left(x^{-\alpha+1/2}
    \calM_{\alpha+1}^1(y^{\alpha+3/2}f,x)-x^{-\alpha+1/2}
    \calM_{\alpha+2n+3}^1(y^{\alpha+3/2}f,x)\right)
\end{multline*}
with $d_n^{-1} = \sum_{k=0}^n \lambda_k = 4(n+1)(\alpha+n+2)$.
\end{propo}
%%%%%%%%%%%%%%%%%%%%%%%%%%%%%%%%%

%%%%%%%%%%%%%%%%%%%%%%%%%%%%%%%%%%%%%%%%%%%%%%
\section{Proof of Theorem~\ref{th:main}}
%\setcounter{equation}{0}
\label{sec:proofmain}
%%%%%%%%%%%%%%%%%%%%%%%%%%%%%%%%%%%%%%%%%%%%%%

The main idea of the proof is the following.
The first summand in the decomposition of $R_n^\alpha(f,x)$,
$x^{-\alpha-1/2}\calM_{\alpha}^0(y^{\alpha+1/2}f,x)$,
may be easily treated since we know appropriate weighted bounds for the
operator $\calM_{\alpha}^0$. Thus we are immediately left with the
operators
\begin{equation*}
  Q_n^\alpha(f,x)
  = R_n^\alpha(f,x) - x^{-\alpha-1/2}
    \calM_{\alpha}^0(y^{\alpha+1/2}f,x).
\end{equation*}
The relevant uniform inequalities for them,
\begin{equation*}
  \|Q_n^\alpha f\|_{L^p(\muag)}
  \le C \|f\|_{L^p(\muag)},\qquad n=1,2,\dots,
\end{equation*}
will be obtained by a proper splitting of $(0,\infty)\times(0,\infty)$ and
then by treating each resulting region separately. It is important to stress
here the fact that in the region that corresponds to large $x$ and $y$ we will
go back to the original definition
of $R_n^\alpha$ as a Ces\`aro type mean of $\{S^\alpha_k\}$.
In consequence, the results of Section~\ref{sec:kernel}
will allow us to write $Q_n^\alpha f$ as a mean of
$x^{-\alpha-1/2} \calM_{\alpha+2k+2}^0(y^{\alpha+1/2}f,x)$,
$k = 0,1,\dots,n$,
and then to treat each resulting term in the aforementioned
region separately.

The weighted bound for $\calM_{\alpha}^0$ is contained in
the following (see, with different degree of generality, \cite{Her},
\cite{CiGPV-DIE}, \cite{St03} and, in particular, \cite[\S\,6.4]{CiGPV}).

%%%%%%%%%%%%%%%%%%%%%%%%%%%%%%%%%
\begin{propo}
\label{pro:Mab}
Let $\alpha>-1$, $\gamma\in \mathbb{R}$ and $1 < p < \infty$. Then
\begin{equation*}
  \| x^{-\alpha-1/2}\calM_{\alpha}^0(y^{\alpha+1/2}f,x)\|_{L^p(\muag)}
  \le C \| f(x) \|_{L^p(\muag)}
\end{equation*}
if and only if~\eqref{eq:condition1} and~\eqref{eq:condition2} are satisfied.
\end{propo}
%%%%%%%%%%%%%%%%%%%%%%%%%%%%%%%%%

Let $\nu=\nu(n)=\alpha+2n+3$, $n=0,1,\dots$, and
$A=(0,4\nu)$, $B=(2\nu,\infty)$. Since
\begin{equation*}
  \chi_{(0,\infty)\times(0,\infty)}
  = \chi_{A\times A} + \chi_{A^c\times B^c}
  + \chi_{B^c\times A^c} + \chi_{B\times B}
  - \chi_{(A\cap B)\times (A\cap B)},
\end{equation*}
by using the notation $Q^\alpha_{n,T,S}f=\chi_T Q^\alpha_n(\chi_S f)$
we have
\begin{multline*}
  \|Q^\alpha_nf\|_{L^p(\muag)}
  \le \|Q^\alpha_{n,A,A}f\|_{L^p(\muag)}
  + \|Q^\alpha_{n,A^c,B^c}f\|_{L^p(\muag)} \\
  + \|Q^\alpha_{n,B^c,A^c}f\|_{L^p(\muag)}
  + \|Q^\alpha_{n,B,B}f\|_{L^p(\muag)}
  + \|Q^\alpha_{n,A\cap B,A\cap B}f\|_{L^p(\muag)}.
\end{multline*}
We will bound each of the five quantities above by $C\|f\|_{L^p(\muag)}$
with $C$ independent of~$n$ and~$f$.
In the rest of this section we write $\|\cdot\|_p$ for the unweighted
$L^p$-norm on~$(0,\infty)$.

%%%%%%%%%%%%%%%%%%%%%%
\subsection{Estimating $\|Q^\alpha_{n,A,A}f\|_{L^p(\muag)}$}
%%%%%%%%%%%%%%%%%%%%%%

Noting that $d_n\sim \nu(n)^{-2}$, it is enough to show the uniform bounds:
\begin{equation*}
  \nu^{-2} \|x^{-\alpha+1/2}\chi_A(x)
    \calM_{\nu}^1(y^{\alpha+3/2}\chi_Af,x)\|_{L^p(\muag)}
  \le C\|\chi_Af\|_{L^p(\muag)}, \quad \nu\ge\alpha+3,
\end{equation*}
and the analogous bound with $\calM_{\nu}^1$ frozen to
$\calM_{\alpha+1}^1$ (but with $A$ associated to $\nu$,
not to $\alpha+1$\,!). In both cases,
taking $g(y)=y^{\alpha+3/2}f(y)$ and using the notation
\begin{equation}
\label{eq:aaa}
  a = -\alpha-1/2+(2(\alpha+\gamma)+1)/p,
\end{equation}
proving the former bounds reduces to showing that
\begin{equation*}
  \nu^{-2}\|x\chi_A(x)\calM_\nu^1(\chi_Ag,x)x^a\|_p
  \le C\|y^{-1}\chi_Agy^a\|_p,
  \quad \nu\ge\alpha+3,
\end{equation*}
and the analogous bound with $\calM_{\nu}^1$ replaced by
$\calM_{\alpha+1}^1$.
These bounds, due to $x\le4\nu$ and $(4\nu)^{-1}\le y^{-1}$,
are further reduced to
\begin{equation}
\label{eq:red}
  \|\chi_A(x)\calM_\nu^1(\chi_Ag,x)x^a\|_p
  \le C \|\chi_Agy^a\|_p,
  \quad \nu\ge\alpha+3,
\end{equation}
and the one with $\calM_{\nu}^1$ replaced by
$\calM_{\alpha+1}^1$.

In~\cite{CiStV}, the following result on the uniform boundedness
of the Bochner-Riesz operators for the Hankel
transform was found (actually, in a more general setting).

%%%%%%%%%%%%%%%%%%%%%%%%%%%%%
\begin{propo}[{\cite[Theorem~1]{CiStV}}]
\label{pro:acotSnu}
Let $\beta>-1$ and $1<p<\infty$.
The Bochner-Riesz multipliers of order one for $\calH_\nu$
given by~\eqref{eq:B-Rd}, satisfy the uniform estimate
\begin{equation*}
  \| \calM_{\nu}^1(g,x) x^a \|_p \le C \| g(x) x^a \|_p,
  \qquad \nu\ge \beta,
\end{equation*}
if and only if $-1/p - (\beta+1/2) < a < 1 - 1/p + (\beta+1/2)$
and $-1/p - 1 < a < 2 - 1/p$.
\end{propo}
%%%%%%%%%%%%%%%%%%%%%%%%%%%%%%%%%%%%%%%%%%%

Thus, taking $\beta=\alpha+1$ gives~\eqref{eq:red}
and its companion once we assume that
\begin{equation*}
  \max\left\{1,\frac{4(\alpha+\gamma+1)}{2\alpha+5}\right\}
  < p <
  \begin{cases}\frac{4(\alpha+\gamma+1)}{2\alpha-1},
    &\text{ if $\alpha>1/2$},\\
    \infty,&\text{ if $-1<\alpha\le 1/2$},
  \end{cases}
\end{equation*}
and
\begin{equation*}
%\label{eq:hypAA''}
  0 < \alpha+\gamma+1 < (\alpha+3/2)p
\end{equation*}
are satisfied. Finally, note that the two conditions
are implied by~\eqref{eq:condition1} and~\eqref{eq:condition2}.

%%%%%%%%%%%%%%%%%%%%%%
\subsection{Estimating $\|Q^\alpha_{n,B^c,A^c}f\|_{L^p(\muag)}$ and
$\|Q^\alpha_{n,A^c,B^c}f\|_{L^p(\muag)}$}
%%%%%%%%%%%%%%%%%%%%%%

Applying an argument from the previous subsection and
taking $g(y)=y^{\alpha+1/2}f(y)$ it is enough to
prove the uniform inequalities
\begin{equation}
\label{eq:one}
  \nu^{-2} \|x\chi_{B^c}(x)\calM_\nu^1(y\chi_{A^c}g,x)x^{a}\|_p
  \le C \|g(y)y^a\|_p, \quad \nu\ge\alpha+3,
\end{equation}
and
\begin{equation}
\label{eq:two}
  \nu^{-2} \|x\chi_{A^c}(x)\calM_\nu^1(y\chi_{B^c}g,x)x^{a}\|_p
  \le C \|g(y)y^a\|_p, \quad \nu\ge\alpha+3,
\end{equation}
with $a$ given by~\eqref{eq:aaa}, and its counterparts with
$\calM_{\nu}^1$ replaced by $\calM_{\alpha+1}^1$.
This will be done in Propositions~\ref{pro:BcAc} and~\ref{pro:AcBc}
below by applying uniform pointwise estimates of Bessel
functions which were already used in~\cite{CiStV}: with a constant~$C$,
independent of $\nu$ and~$x$,
\begin{equation}
\label{eq:Phi}
  |x^{1/2}J_\nu(x)| \le C \Phi_\nu(x)
\end{equation}
where, for $\nu \ge 1$,
\begin{equation*}
  \Phi_\nu(x) =
  \begin{cases}
    \frac{x^{\nu+1/2}}{\Gamma(\nu+1)2^\nu},& 0<x<\nu/2,\\
    \nu^{1/4}(\nu^{1/3}+|x-\nu|)^{-1/4},& \nu/2<x<2\nu,\\
    1,& 2\nu<x<\infty,
  \end{cases}
\end{equation*}
and, for $-1 < \nu < 1$,
\begin{equation*}
  \Phi_\nu(x) =
  \begin{cases}
    \frac{x^{\nu+1/2}}{\Gamma(\nu+1)},& 0<x<1/2,\\
    1,& 1/2<x<\infty.
  \end{cases}
\end{equation*}
Moreover, there exists a constant $C$ independent of $\nu>-1$, such
that, for $(x,y) \in B^c \times A^c$ or $(x,y) \in A^c \times B^c$,
we have
\begin{equation}
\label{eq:acotcalK}
  |\calM_\nu^1(x,y)|
  \le C \frac{\Phi_\nu(x)\Phi_\nu(y)}{|x-y|^2};
\end{equation}
see~\cite[Lemma~4]{CiStV} and a slight modification
of~\cite[Lemma~2]{CiStV}.

%%%%%%%%%%%%%%%%%%%%%%%%%%%%%%%%%
\begin{propo}
\label{pro:BcAc}
Let $\alpha>-1$, $\alpha+\gamma>-1$, $1<p<\infty$, and $a$ be given
by~\eqref{eq:aaa}. If $p<p_1(\alpha,\gamma)$,
then the inequality~\eqref{eq:one} and its aforementioned counterpart
are satisfied.
\end{propo}
%%%%%%%%%%%%%%%%%%%%%%%%%%%%%%%%%

\begin{proof}
We consider only the case of~\eqref{eq:one}.
The proof of its mutant, with $\calM_{\alpha+1}^1$
in place of $\calM_\nu^1$,
requires minor modifications only.
For $(x,y) \in B^c \times A^c$, $x \le 2\nu$. Thus
\begin{equation*}
  \left( \int_0^{2\nu} \Big| \frac{x}{\nu^2}
    \calM_\nu^1(yg\chi_{A^c},x)x^a \Big|^p \,dx \right)^{1/p}
  \le \frac{C}{\nu} \left(\int_0^{2\nu} \big|
    \calM_\nu^1(yg\chi_{A^c},x)x^a \big|^p \,dx\right)^{1/p}.
\end{equation*}
Moreover, in that region,  $y-x\ge y/2$, so $1/|x-y|^2 \le 4/y^2$.
Then, by using~\eqref{eq:acotcalK} and H\"older's inequality,
\begin{align*}
  \big| \chi_{B^c}(x) \calM_\nu^1(yg\chi_{A^c},x)x^a \big|
  &\le C \chi_{B^c}(x)x^a\Phi_\nu(x)
    \left(\int_0^\infty\chi_{A^c}(y) y^{-1}|g(y)| \Phi_\nu(y) \,dy \right) \\
  &\le C \chi_{B^c}(x) x^a \Phi_\nu(x)\,\| g(y) y^a\|_p \,
    \|\chi_{A^c}(y) y^{-a-1} \Phi_\nu(y) \|_{p'}
\end{align*}
and thus
\begin{multline*}
 \nu^{-2} \| \chi_{B^c}(x) \calM_\nu^1(yg\chi_{A^c},x)x^a \|_p \\
  \le C \nu^{-1}\, \| g(y) y^a\|_p\,
  \|\chi_{A^c}(y) y^{-a-1} \Phi_\nu(y) \|_{p'}\,
  \|\chi_{B^c}(x) x^a \Phi_{\nu}(x) \|_p.
\end{multline*}
Therefore, to establish the proposition we only need to check that
the product of the last two terms is bounded by $C\nu$.
This easily follows, by a direct calculation, once we have $(-a-1)p'<-1$
(i.e., $p<p_1(\alpha,\gamma)$),
from the definition of~$\Phi_\nu$.
\end{proof}

%%%%%%%%%%%%%%%%%%%%%%%%%%%%%%%%%
\begin{propo}
\label{pro:AcBc}
Let $\alpha>-1$, $\alpha+\gamma>-1$, $1<p<\infty$ and $a$ be given
by~\eqref{eq:aaa}.
If $p_0(\alpha,\gamma)<p$, then the inequality~\eqref{eq:two}
and its aforementioned counterpart are satisfied.
\end{propo}
%%%%%%%%%%%%%%%%%%%%%%%%%%%%%%%%%

\begin{proof}
The proof is very similar to that of the previous proposition.
Moreover, a duality argument can also  be used, because the operators
from Propositions~\ref{pro:BcAc} and~\ref{pro:AcBc} are adjoint in
the following sense:
\begin{equation*}
  \langle  x\chi_{B^c} \calM_\nu^1(y^{-a}yg_1\chi_{A^c},x) x^a, h_1 \rangle
  = \langle  y\chi_{A^c} \calM_\nu^1(x^axh_1\chi_{B^c},y) y^{-a}, g_1 \rangle.
\end{equation*}
Thus, the hypothesis $(-a-1)p' < -1$ from the proof of Proposition~\ref{pro:BcAc}
becomes $(a-1)p < -1$, i.e.,~$p_0(\alpha,\gamma)<p$.
\end{proof}

%%%%%%%%%%%%%%%%%%%%%%
\subsection{Estimating $\|Q^\alpha_{n,B,B}f\|_{L^p(\muag)}$
and $\|Q^\alpha_{n,A\cap B, A\cap B}f\|_{L^p(\muag)}$}
%%%%%%%%%%%%%%%%%%%%%%

Start with the case of $B\times B$. Here we do not use the
decomposition of $Q_n^\alpha(f,x)$ that follows from
Proposition~\ref{pro:3sum}. Instead, we apply
earlier results of Section~\ref{sec:kernel} that lead
to writing $Q_n^\alpha(f,x)$ as the mean
\begin{equation*}
  Q_n^\alpha(f,x) = -\frac{d_n}{2} \sum_{k=0}^n
  \lambda_k x^{-\alpha-1/2}\calM_{\alpha+2k+2}^0(y^{\alpha+1/2}f,x).
\end{equation*}
The required bound
\begin{equation*}
  \|Q^\alpha_{n,B,B}f\|_{L^p(\muag)} \le C\|f\|_{L^p(\muag)}
\end{equation*}
then follows from the following.

%%%%%%%%%%%%%%%%%%%%%%%%
\begin{propo}
\label{pro:BB}
Let $\alpha>-1$, $\alpha+\gamma>-1$, $1<p<\infty$ and,
for simplicity, write $\kappa=\alpha+2k+2$.
If the assumption~\eqref{eq:condition2} is satisfied then
\begin{equation*}
  \| x^{-\alpha-1/2}\chi_{B}(x)
    \calM_{\kappa}^0(y^{\alpha+1/2}\chi_{B}f,x) \|_{L^p(\muag)}
  \le C \| f(x) \|_{L^p(\muag)},
  \quad 0 \le k \le n,
\end{equation*}
with a constant $C$ independent of~$n$.
\end{propo}
%%%%%%%%%%%%%%%%%%%%%%%%%

\begin{proof}
Denoting by $H(g,x)$ the Hilbert transform of~$g$,
\begin{equation*}
  H(g,x) = \text{P.V.} \int_0^\infty \frac {g(y)}{x-y}\,dy,
\end{equation*}
and performing the change of variables
$x \mapsto x^{1/2}$ and $y \mapsto y^{1/2}$ in~\eqref{eq:phi3},
it is then enough to prove
\begin{multline}
\label{eq:JJ'}
  \int_0^\infty \big| \chi_B(\sqrt{x}) x^{-\alpha/2} J_\kappa(\sqrt{x})
    H(\chi_B(\sqrt{y}) f(\sqrt{y}) y^{(1+\alpha)/2} J'_\kappa(\sqrt{y}), x)
    \big|^p x^{\alpha+\gamma}\,dx \\
  \le C \int_0^\infty\big| f(\sqrt{x})\big|^p x^{\alpha+\gamma}\,dx
\end{multline}
and the analogous inequality with $J_\kappa(\sqrt{x})$ replaced by
$J'_\kappa(\sqrt{x})$ and $J'_\kappa(\sqrt{y})$ replaced by $J_\kappa(\sqrt{y})$.
We treat~\eqref{eq:JJ'} only, the proof of its companion is completely analogous.
By~\eqref{eq:Phi} (and similar bounds for $J'_\nu$, see \cite{CiStV}), it
is clear that $|J_\kappa(x)| \le Cx^{-1/2}$ and
$|J'_\kappa(x)| \le Cx^{-1/2}$ for $x \ge 2\nu$
with a constant $C$ independent of $x$, $\kappa$ and~$\nu$ ($\nu\ge\alpha+2>1$).

We will use the fact that the Hilbert
transform is a bounded operator on weighted $L^p$ spaces with weights from
$A_p$. It is clear that $x^{\alpha+\gamma \pm p/4 - \alpha p/2} \in A_p$
if and only if $-1 < \alpha+\gamma \pm p/4 - \alpha p/2 < p-1$.
Here, the relevant inequalities are $-1 < \alpha+\gamma - p/4 - \alpha p/2$
and $ \alpha+\gamma + p/4 - \alpha p/2 < p-1$,
which are equivalent to~\eqref{eq:condition2} when  considered together.
Take
\begin{equation*}
  g(y) = \chi_B(\sqrt{y}) f(\sqrt{y}) y^{(1+\alpha)/2} J'_\kappa(\sqrt{y}).
\end{equation*}

To prove~\eqref{eq:JJ'} we apply successively
$|J_\kappa(\sqrt{x})| \le Cx^{-1/4}$, $x^{\alpha+\gamma-p/4-\alpha
p/2} \in A_p$, and $|J'_\kappa(\sqrt{x})| \le Cx^{-1/4}$,
obtaining
\begin{align*}
  \int_0^\infty \big| \chi_B(\sqrt{x}) x^{-\alpha/2} J_\kappa(\sqrt{x})
    H(g, x) \big|^p\,x^{\alpha+\gamma}\,dx
  &\le C \int_0^\infty  |H(g, x)|^p x^{\alpha+\gamma-p/4-\alpha p/2}\,dx\\
  &\le C \int_0^\infty  |g(x)|^p x^{\alpha+\gamma-p/4-\alpha p/2}\,dx \\
  &\le  C \int_0^\infty \big|f(\sqrt{x})\big|^p x^{\alpha+\gamma}\,dx.
\end{align*}
\end{proof}

The bound for $\|Q^\alpha_{n,A\cap B, A\cap B}f\|_{L^p(\muag)}$ follows
from the uniform estimate
\begin{equation*}
  \| x^{-\alpha-1/2}\chi_{A\cap B}(x)\calM_{\kappa}^0(y^{\alpha+1/2}
  \chi_{A\cap B}f,x) \|_{L^p(\muag)}
  \le C \| f(x) \|_{L^p(\muag)},
\end{equation*}
$0 \le k \le n$, with $C$ independent of $n$. This can be obtained  by applying
Proposition~\ref{pro:BB} to the function $f(y)\chi_{A\cap B}(y)$ and
taking into account that
\begin{equation*}
  |x^{-\alpha-1/2}\chi_{A\cap B}(x)\calM_{\kappa}^0(y^{\alpha+1/2}
    \chi_{A\cap B}f,x)|
  \le | x^{-\alpha-1/2}\chi_{B}(x)\calM_{\kappa}^0(y^{\alpha+1/2}
    \chi_{A\cap B}f,x)|.
\end{equation*}
The proof of Theorem~\ref{th:main} is completed.

%%%%%%%%%%%%%%%%%%%%%%%%%%%%%%%%%%%%%%%%%%%%%%
\section{Proof of Theorem~\ref{th:Bp=Ep}}
%\setcounter{equation}{0}
\label{sec:Bp=Ep}
%%%%%%%%%%%%%%%%%%%%%%%%%%%%%%%%%%%%%%%%%%%%%%

We follow closely the proof of Theorem~4 in \cite{CiGPV}, where we
refer the reader for details. Clearly,
the assumptions $p_{0}(\alpha+\gamma)<p<p_1(\alpha+\gamma)$
and $\alpha+\gamma\ge-1/2$ are needed to have the
space $E_{p,\alpha+\gamma}$ well defined and assure good properties of
it, see a comment in Section~\ref{sec:intro}. Moreover, in our argument the
case $p=2$ is crucial since it is required when considering other
cases, $p<2$ and $p>2$. Therefore, in the range of $p$'s, $p_0(\alpha,\gamma)<p$,
for which the spaces $B_{p,\alpha,\gamma}$ are well defined, $p=2$ must be
included, which forces us to impose the condition
$\gamma<1/2$. Clearly, with the last assumption the spaces
$B_{p,\alpha,\gamma}$ are well defined for all $p\ge 2$.
Thus, apart of $\alpha>-1$, in what follows we assume
the conditions $\alpha+\gamma\ge-1/2$,
$p_{0}(\alpha+\gamma)<p<p_1(\alpha+\gamma)$, and $\gamma<1/2$ to hold.

The formula (which holds for $\alpha>-1$, $\gamma<1$ and $\alpha+\gamma> -1$,
see \cite[Lemma~3]{CiGPV})
\begin{equation*}
  H_{\alpha+\gamma}(j_n^\alpha,x)
  = 2^{1/2+\gamma}
    \frac{\sqrt{\alpha+2n+1}\,\Gamma(n+1)}{\Gamma(n+1-\gamma)}
    (1-x^2)^{-\gamma} P_n^{(\alpha+\gamma,-\gamma)}(1-2x^2)\chi_{[0,1]}(x).
\end{equation*}
shows the inclusion $B_{p,\alpha,\gamma} \subseteq E_{p,\alpha+\gamma}$.
It remains to check that in fact
the inclusion becomes the identity (with the additional assumption
$\alpha+\gamma+1 < \frac{p}4(2\alpha+4\gamma+3)$ if $p<2$).

\textit{Case $p=2$.}
The above formula, the completeness of the Jacobi system and the fact that
$H_{\alpha+\gamma}$ is an automorphism of $L^2(\muag)$ prove the claim,
i.e.~$B_{2,\alpha,\gamma} = E_{2,\alpha+\gamma}$.

\textit{Case $p>2$.} Here the crucial fact is the result of {Case $p=2$} and
continuity and density of the inclusion
$E_{2,\alpha+\gamma} \subset E_{p,\alpha+\gamma}$.

\textit{Case $p<2$.} The required
equality follows by showing that the only
functional $T\in(E_{p,\alpha+\gamma})'$ such that $T(f)=0$ for all
$f\in B_{p,\alpha,\gamma}$ is
$T=0$. Since dual of $E_{p,\alpha+\gamma}$ is naturally
identified with $E_{p',\alpha+\gamma}$,
there exists $g\in E_{p',\alpha+\gamma}$, such that
$T(f) = \int_0^\infty g(x)f(x)\,\muag(x)$, for every
$f\in E_{p,\alpha+\beta}$. With the additional assumption
$\alpha+\gamma+1 < \frac{p}4(2\alpha+4\gamma+3)$, we have
$j_n^{\alpha+\gamma} \in B_{p,\alpha,\gamma}$,
see \cite[Lemma~2]{CiGPV}. Therefore, in particular,
$c_n^{\alpha+\gamma}(g)
= \int_0^\infty g(x)j_n^{\alpha+\gamma}(x)\,\muag(x)=0$,
for $n=0,1,2,\dots$. Now, from the preceding case, we have
$B_{p',\alpha+\gamma,0} = E_{p',\alpha+\gamma}$,
so that $g\in B_{p',\alpha+\gamma,0}$. This
fact, taking into account that with the hypothesis
$p_0(\alpha,\gamma)<p<2$ we can apply Theorem~\ref{th:main},
together with $c_n^{\alpha+\gamma}(g) = 0$ gives
$0 = \lim_{n\to\infty}R_n^{\alpha+\gamma}g = g$
in $L^{p'}(\muag)$.

The proof of Theorem~\ref{th:Bp=Ep} is completed.

%%%%%%%%%%%%%%%%%%%%%%%%%%%%%%%%%%%%%%%%%%%%%%%%%%%%%%
\begin{thebibliography}{00}
%%%%%%%%%%%%%%%%%%%%%%%%%%%%%%%%%%%%%%%%%%%%%%%%%%%%%%

\bibitem{BeSt} J. J. Betancor and K. Stempak,
Relating multipliers and transplantation for Fourier-Bessel expansions and
Hankel transform,
\textit{T\^ohoku Math. J.}
\textbf{53} (2001), 109--129.

\bibitem{Ci} \'O. Ciaurri,
{``Aproximaci\'on de funciones cuya transformada de Hankel est\'a
soportada en el intervalo $[0,1]$''},
Doctoral dissertation,
Universidad de La Rioja, Logro\~no, Spain, 2000.

\bibitem{CiGPV} \'O. Ciaurri, J. J. Guadalupe, M. P\'erez, and J. L. Varona,
Mean and almost everywhere convergence of Fourier-Neumann series,
\textit{J. Math. Anal. Appl.}
\textbf{236} (1999), 125--147.

\bibitem{CiGPV-DIE} \'O. Ciaurri, J. J. Guadalupe, M. P\'erez, and J. L. Varona,
Solving dual integral equations on Lebesgue spaces,
\textit{Studia Math.}
\textbf{142} (2000), 253--267.

\bibitem{CiV} \'O. Ciaurri and J. L. Varona,
A uniform boundedness for Bochner-Riesz operators related
to the Hankel transform,
\textit{J. Inequal. Appl.} \textbf{7} (2002), 759--777.

\bibitem{CiStV} \'O. Ciaurri, K. Stempak, and J. L. Varona,
Uniform two-weight inequalities for Hankel transform Bochner-Riesz means,
preprint, December 2002.

\bibitem{Dur} A. J. Dur\'an,
On Hankel transform,
\textit{Proc. Amer. Math. Soc.}
\textbf{110} (1990), 417--424.

\bibitem{Har} G. H. Hardy,
{``Divergent Series''},
Oxford at the Clarendon Press, Oxford, 1949.

\bibitem{Her} C. S. Herz,
On the mean inversion of Fourier and Hankel transforms,
\textit{Proc. Nat. Acad. Sci. U.S.A.}
\textbf{40} (1954), 996--999.

\bibitem{Kno} K. Knopp,
{``Theory and Application of Infinite Series''},
Hafner, New York, 1971.

\bibitem{Po} E. L. Poiani,
Mean Ces\`aro summability of Laguerre and Hermite series,
\textit{Trans. Amer. Math. Soc.}
\textbf{173} (1972), 1--31.

\bibitem{St91} K. Stempak,
Almost everywhere summability of Laguerre series,
\textit{Studia Math.}
\textbf{100} (1991), 129--147.

\bibitem{St03} K. Stempak,
Uniform two-weight norm inequalities for Hankel transform partial sum operators,
\textit{Houston J. Math.}
\textbf{29} (2003), 1045--1063.

\bibitem{V} J. L. Varona,
Fourier series of functions whose Hankel transform is supported on $[0,1]$,
\textit{Constr. Approx.}
\textbf{10} (1994), 65--75.

\end{thebibliography}

\end{document}

